GraphRAG integrates (knowledge) graphs with large language models (LLMs) to improve reasoning accuracy and contextual relevance. 
Despite its promising applications and strong relevance to multiple research communities, such as databases and natural language processing, GraphRAG currently lacks modular workflow analysis, systematic solution frameworks, and insightful empirical studies. 
To bridge these gaps, we propose {\bf LEGO-GraphRAG}, a modular framework that enables: {\bf 1)} fine-grained decomposition of the GraphRAG workflow, {\bf 2)} systematic classification of existing techniques and implemented GraphRAG instances, and {\bf 3)} creation of new GraphRAG instances. 
Our framework facilitates comprehensive empirical studies of GraphRAG on large-scale real-world graphs and diverse query sets, revealing insights into balancing reasoning quality, runtime efficiency, and token or GPU cost, that are essential for building advanced GraphRAG systems. 
%The implementation of our framework and the technical report are available at~\ourcode.


%Despite its potential and appeal to diverse research communities, including databases and natural language processing, GraphRAG lacks modular workflow analysis, a systematic framework for {\color{red} organizing} its diverse components, and comprehensive empirical evaluations of their interoperability.
%Based on the proposed framework, we conduct extensive empirical studies on large-scale real-world graphs and diverse GraphRAG query sets, identifying key findings, including: {\bf 1)} GraphRAG must balance quality, efficiency, and cost across three hierarchical levels: modules, method types, and semantic models. {\bf 2)} Extracting query-relevant subgraphs is the primary efficiency bottleneck for GraphRAG in large-scale graphs and remains overlooked. {\bf 3)} Graph structural information enables efficient solutions for GraphRAG, while semantic information is vital for improving complex query quality. The optimal approach uses structural information to identify query-relevant subgraphs and semantic information to retrieve reasoning paths.




\begin{comment}
	内容...

    {\color{red} To address hallucination and knowledge deficiencies in large language models (LLMs), Graph-based Retrieval-Augmented Generation (GraphRAG) has demonstrated significant potential by leveraging graphs as external resources to enhance LLM reasoning capabilities. Nevertheless, existing GraphRAG approaches remain in their nascent stages, and a comprehensive evaluation and analysis are still lacking.
    In this paper, we present the first systematic evaluation and analysis of GraphRAG methodologies. We introduce a modular GraphRAG framework and conduct extensive experiments on widely recognized benchmark datasets, evaluating the performance of GraphRAG methods across various levels. Our findings elucidate the strengths and limitations of current approaches, offering actionable insights for advancing future GraphRAG systems.}


    {\color{blue}
    1. Briefly introduce GraphRAG, GraphRAG is an emerging field of research

    2. Two problems: 

    1) Non-uniform evaluation framework

    2) Non-modular analysis framework


    {\bf Non-Unified Evaluation Framework:} Recent research on GraphRAG is rather heterogeneous and has not been evaluated under a unified framework, making it difficult to clarify the actual utility of the proposed key algorithms and components for GraphRAG.

    {\bf Non-modular analysis framework:} There is no modular framework to disassemble the execution process of GraphRAG in a fine-grained way, so as to easily analyze the performance trade-offs of GraphRAG in different scenarios.
    3. We do xxxxx
    }
    
      To address hallucination and information deficiency in Large Language Models (LLMs), Retrieval-Augmented Generation (RAG) has been introduced as an effective approach. By querying relevant knowledge during the retrieval process and leveraging the knowledge during generation, RAG helps mitigate hallucination and enhances the model's access to domain-specific information. 
    GraphRAG is an emerging field of research that enhances the reasoning capabilities of LLMs by utilizing graphs instead of text documents as external resources. However, existing GraphRAG methods lack comprehensive evaluation and analysis, facing two primary challenges. One is the {\bf Non-Unified Evaluation Framework}. Recent research on GraphRAG is heterogeneous and has not been assessed under a unified framework, which complicates the clarification of the actual utility of the proposed key algorithms and components. 
    The other is the {\bf Non-Modular Analysis Framework}. There is an absence of a modular framework to decompose the execution process of GraphRAG in a fine-grained manner, hindering the analysis of performance trade-offs in various scenarios. 
    To resolve these issues, we propose the first systematic evaluation and analysis framework of GraphRAG methodologies. We introduce a modular GraphRAG framework and conduct extensive experiments on widely recognized benchmark datasets, assessing the performance of GraphRAG methods at different levels. Our findings elucidate the strengths and limitations of current approaches, providing actionable insights for future GraphRAG systems.
\end{comment}




